\documentclass[12pt]{article}
\usepackage[T1]{fontenc}
\usepackage[utf8]{inputenc}
\usepackage{lmodern}
\usepackage{textcomp}
\usepackage{enumitem}
\usepackage{geometry}
\geometry{margin=1in}
\begin{document}
\begin{center}
Answer Key - Psychology - Graduate Level Assessment 14 \
\small (Answers are ordered exactly as in the question paper.)
\end{center}

\section*{Multiple Choice}
\begin{enumerate}[label=\arabic*.]
\item \textbf{Answer:} D
\\ \textit{Options:} A) when they are not on any form of medication; B) just before beginning antidepressant medication; C) during the first few days after starting antidepressant medication; D) as they begin to recover from their depression; E) during periods of remission from depression
\\ \textit{Note:} As individuals begin to recover from depression, they may experience a surge of energy or clarity that can lead them to act on suicidal thoughts they had been previously too lethargic to pursue, making them more dangerous to themselves during this critical period.
\item \textbf{Answer:} A
\\ \textit{Options:} A) Three subtypes of anxiety disorders are listed in DSM-III; B) Five subtypes of dissociative disorders are listed in DSM-III; C) Six subtypes of depression disorders are listed in DSM-III; D) Eight subtypes of somatoform disorders are listed in DSM-III; E) Seven subtypes of eating disorders are listed in DSM-III
\\ \textit{Note:} The correct answer is accurate because DSM-III categorizes anxiety disorders into three distinct subtypes: phobic disorders, panic disorder, and generalized anxiety disorder, each with specific diagnostic criteria and characteristics.
\item \textbf{Answer:} A
\\ \textit{Options:} A) parietal lobe; B) occipital lobe; C) hypothalamus; D) cerebellum; E) hippocampus
\\ \textit{Note:} The parietal lobe is involved in sensory perception and integration, and damage to this area in Phineas Gage was crucial in understanding the effects of brain injury on personality and behavior, making it the correct answer in the context of his accident.
\item \textbf{Answer:} A
\\ \textit{Options:} A) interneuron, afferent neuron, efferent neuron, receptor, effector; B) effector, sensory neuron, afferent neuron, interneuron, receptor; C) sensory neuron, interneuron, afferent neuron, efferent neuron, effector; D) sensory neuron, afferent neuron, interneuron, efferent neuron, receptor; E) receptor, efferent neuron, interneuron, afferent neuron, affector
\\ \textit{Note:} The correct answer represents the sequence of neuronal communication in a reflex arc, where the afferent neuron conveys sensory information to the interneuron in the spinal cord, which then relays the impulse to the efferent neuron that activates the effector, thereby facilitating a quick, involuntary response to a stimulus.
\item \textbf{Answer:} A
\\ \textit{Options:} A) secondary intervention.; B) quaternary prevention.; C) secondary prevention.; D) psychological intervention.; E) advanced prevention.
\\ \textit{Note:} Providing education and job training to adolescents and young adults recently released from a drug treatment program is considered secondary intervention because it aims to prevent relapse and support recovery by addressing the social and economic factors that contribute to substance abuse.
\end{enumerate}

\section*{True / False}
\begin{enumerate}[label=\arabic*.]
\setcounter{enumi}{5}
\item \textbf{Answer:} False
\\ \textit{Options:} A) True; B) False
\\ \textit{Note:} Thermoreceptors primarily transduce thermal energy into electrochemical energy, rather than mechanical energy, to facilitate the sensation of temperature.
\item \textbf{Answer:} False
\\ \textit{Options:} A) True; B) False
\\ \textit{Note:} Cognitive therapy primarily emphasizes the modification of maladaptive thought patterns rather than relying on nonverbal cues, which are more central to approaches like humanistic or psychodynamic therapies.
\item \textbf{Answer:} True
\\ \textit{Options:} A) True; B) False
\\ \textit{Note:} Harlow's study with rhesus monkeys demonstrated that while emotional bonds are crucial for development, the specific attachment styles observed in humans may involve more complex social and cognitive factors beyond early emotional bonding.
\item \textbf{Answer:} True
\\ \textit{Options:} A) True; B) False
\\ \textit{Note:} Job satisfaction is shaped by a complex interplay of individual factors, such as personal values and life circumstances, as well as external elements like workplace environment and organizational culture, highlighting the multifaceted nature of how satisfaction at work is achieved.
\item \textbf{Answer:} False
\\ \textit{Options:} A) True; B) False
\\ \textit{Note:} Problem-focused coping involves utilizing external resources and strategies to address and change the stressor itself, rather than relying solely on internal resources.
\end{enumerate}

\section*{Long Form Response}
\begin{enumerate}[label=\arabic*.]
\setcounter{enumi}{10}
\item \textbf{Answer:} Nutritional content, food storage methods, and the timing of meals play significant roles in shaping taste perception. The nutritional composition of food influences not only its flavor profile but also the physiological state of the consumer; for instance, a diet high in sugars can alter taste sensitivity, often leading to preferences for sweeter foods. Food storage methods can affect taste by causing chemical changes, such as the degradation of flavor compounds in improperly stored items, which diminishes taste intensity. Furthermore, the timing of meals, particularly the duration since the last meal, affects hunger levels and sensory sensitivity; when hungry, individuals may experience heightened taste perception, enhancing the enjoyment of flavors. Together, these factors underscore the complex interplay between physiological states and sensory experiences, illustrating how external and internal conditions influence our taste perceptions.
\\ \textit{Note:} correct because it effectively connects the nutritional composition, storage methods, and timing of meals to their collective impact on taste perception, emphasizing how these factors can alter flavor profiles, taste sensitivity, and overall sensory experience.
\item \textbf{Answer:} Research findings indicate that gender differences in conversation style often manifest in women asking more questions than men during interactions. This tendency is attributed to women's socialization, which emphasizes relational communication and active listening, thereby fostering an environment that encourages inquiry and engagement. Women are generally more attuned to the dynamics of conversation and use questions to promote inclusivity and understanding. Additionally, studies suggest that women may use questioning as a strategy to build rapport and validate the perspectives of their conversational partners, reflecting a collaborative approach to dialogue. In contrast, men may adopt a more assertive communication style, often prioritizing information delivery over relational dynamics, resulting in fewer questions.
\\ \textit{Note:} correct because it highlights how women's socialization fosters relational communication and active listening, leading them to ask more questions to enhance inclusivity and understanding in conversations.
\end{enumerate}

\vfill
\noindent\textit{End of Answer Key}
\end{document}